\begin{tikzpicture}

\node[box,text width=3.4cm] (a) at (0,0) {输入分区\\读取 offset 0\\下一待提交位移为 1};
\node[good,text width=4.4cm] (t) at (5.2,0) {同一 Kafka 事务\\写输出记录\\写消费组 offset=1\\提交或中止决定};
\node[box,text width=3.7cm] (c) at (11,0) {read\_committed\\提交：看见输出\\中止：过滤输出};
\draw[arr] (a)--(t);\draw[arr] (t)--(c);
\node[warn,text width=4.4cm] (x) at (5.2,-2.7) {消费数据库／外部付款\\各自拥有事务和原身份\\不在该 Kafka 提交中};
\draw[ctl] (a.south)--(x.west);
\node[note,text width=12cm] at (5,-4.4) {事务中止后，应用把本地读取位置恢复到已提交位移。过滤中止记录与重新处理原输入共同组成恢复。};

\end{tikzpicture}
